% =====================================================================
%  第 3 章：格点上的纯规范理论
%  Gattringer & Lang, "Quantum Chromodynamics on the Lattice" (LNP 788)
%  Reconstructed from extract/ch03.txt; equations auto-numbered.
% =====================================================================

\chapter{格点上的纯规范理论}

在上一章中，我们开始了在格点上构造 QCD 的工作。我们讨论了 QCD 的一个特别重要的特征，即胶子的自相互作用。这种自相互作用使得纯胶动力学（即不含夸克的 QCD）成为一个有趣且高度非平凡的理论。纯胶动力学展现出颜色禁闭，这是完整 QCD 的一个重要性质。由于胶动力学比包含夸克的 QCD 容易处理得多，它是研究禁闭及其内在机制的重要课题。

在本章中，我们将通过讨论链接变量的积分测度来完成格点上胶动力学的构造。在关于规范不变性和规范固定的一节之后，我们讨论两个静态夸克源之间的势。静态夸克势随后可以用来设定标度，即确定格点间距 $a$ 作为规范耦合 $g$ 的函数的值。

在进入这些主题之前，让我们简要汇集定义格点上纯规范理论的公式。可观测量 $O$ 的期望值（其中 $O$ 也可以是若干项的乘积）由下式给出
\begin{equation}
  \ev{O} = \frac{1}{Z}\int D[U]\,e^{-S_G[U]}\,O[U],
\end{equation}
其中配分函数 $Z$ 定义为
\begin{equation}
  Z = \int D[U]\,e^{-S_G[U]}.
\end{equation}
链接变量的积分测度是乘积测度
\begin{equation}
  D[U] = \prod_{n\in\Lambda}\prod_{\mu=1}^{4} dU_\mu(n).
\end{equation}
对于规范场作用量，我们引入了威尔逊作用量
\begin{equation}
  S_G[U] = \frac{\beta}{3}\sum_{n\in\Lambda}\sum_{\mu<\nu}
           \Re\,\tr\!\left[1 - U_{\mu\nu}(n)\right],
\end{equation}
但我们再次指出，其他表述形式也同样在使用。在 (3.4) 中我们使用一个广泛使用的缩写
\begin{equation}
  \beta = \frac{6}{g^2},
\end{equation}
即所谓的逆耦合。注意，这是 $\beta$ 在 SU(3) 情形下的定义。一般情形 SU($N$) 的表达式在 (3.93) 中给出。

\section{Haar 测度}

在 QCD 格点作用量的构造中，链接变量 $U_\mu(n)$ 作为 SU(3) 的元素被引入。我们现在讨论 (3.3) 的乘积测度中各个测度 $dU_\mu(n)$ 的定义，该乘积测度将链接变量 $U_\mu(n)$ 在整个 SU(3) 群流形上积分。我们以一种更一般的方式来考虑这个问题，讨论测度 $dU$，其中 $U$ 是一个李群的元素，即一个其元素连续依赖于某些参数的群的元素。

\subsection{规范场测度与规范不变性}

对我们规范场测度的一个重要限制来自我们理论的规范不变性。在上一章中，我们把链接变量的规范变换引入为
\begin{equation}
  U_\mu(n) \to U'_\mu(n) = \Omega(n)\,U_\mu(n)\,\Omega(n+\hat\mu)^\dagger .
\end{equation}
群值矩阵 $\Omega(n)$ 可以在每个格点 $n$ 处独立地选取。规范场的作用量 $S_G[U]$ 在这些变换下是不变的，我们有
\begin{equation}
  S_G[U'] = S_G[U].
\end{equation}
与任何积分一样，路径积分的结果应当在变量变换下保持不变，特别是应在规范变换 (3.6) 下保持不变。对于配分函数，这一要求为
\begin{equation}
  Z = \int D[U]\,e^{-S_G[U]}
    = \int D[U']\,e^{-S_G[U]}
    = \int D[U']\,e^{-S_G[U']},
\end{equation}
其中在最后一步我们使用了规范作用量的规范不变性 (3.7)。比较 (3.8) 中第一个和最后一个表达式，我们得到
\begin{equation}
  D[U] = D[U'].
\end{equation}
利用 $D[U]$ 是乘积测度这一事实（比较 (3.3)），我们导出条件
\begin{equation}
  dU_\mu(n) = dU'_\mu(n) = d\!\left(\Omega(n)\,U_\mu(n)\,\Omega(n+\hat\mu)^\dagger\right),
\end{equation}
用于对各个链接变量进行积分。这一性质是 Haar 测度（哈尔测度）的定义性性质之一。

\subsection{群积分测度}

规范不变性的要求把我们引向数学文献中一个众所周知的结构，即所谓的 Haar 测度。这是用于在连续紧致群 $G$ 上积分的测度。我们已经对单个群变量 $U_\mu(n)$ 的积分导出了条件 (3.10)。由于 $\Omega(n)$ 和 $\Omega(n+\hat\mu)$ 可以独立选取，群元素的测度 $dU$ 必须在与另一个群元素 $V \in G$ 的左乘和右乘下保持不变，即
\begin{equation}
  dU = d(UV) = d(VU).
\end{equation}
最后一个方程，连同归一化条件（对整个 $G$ 的积分）
\begin{equation}
  \int dU\,1 = 1,
\end{equation}
是紧致李群上积分的 Haar 测度的定义性性质。如所预告的那样，我们不仅讨论 SU(3) 的测度，而且讨论更一般的紧致李群 $G$（例如 SU($N$) 或 U(1)）的测度。对于这些情形，我们现在给出 Haar 测度的一个显式构造。设 $U = U(\omega)$ 是指数表示 (A.2) 中紧致李群 $G$ 的一个元素，由一组实数 $\omega^{(k)}$ 参数化。我们已经讨论过（另见附录 A.1），$\partial/\partial\omega^{(k)}\, U(\omega)\, U(\omega)^{-1}$ 位于群的李代数中。基于这些李代数值的对象，我们可以在群上定义度规 $ds^2$ 为（对 $n$、$m$ 使用求和约定）
\begin{equation}
  \begin{split}
    ds^2 &\equiv \tr\!\left[
      \frac{\partial U(\omega)}{\partial \omega^{(n)}} U(\omega)^{-1}
      \left(\frac{\partial U(\omega)}{\partial \omega^{(m)}} U(\omega)^{-1}\right)^{\!\dagger}
    \right] d\omega^{(n)} d\omega^{(m)}\\
    &= g(\omega)_{nm}\, d\omega^{(n)} d\omega^{(m)} .
  \end{split}
\end{equation}
从这个方程我们读出度规张量 $g(\omega)$，
\begin{equation}
  g(\omega)_{nm} = \tr\!\left[
    \frac{\partial U(\omega)}{\partial \omega^{(n)}} U(\omega)^{-1}
    \left(\frac{\partial U(\omega)}{\partial \omega^{(m)}} U(\omega)^{-1}\right)^{\!\dagger}
  \right]
  = \tr\!\left[
    \frac{\partial U(\omega)}{\partial \omega^{(n)}}
    \frac{\partial U(\omega)^{\dagger}}{\partial \omega^{(m)}}
  \right].
\end{equation}
用度规 $g(\omega)$ 我们可以把测度 $dU$ 定义为
\begin{equation}
  dU = c\,\sqrt{\det[g(\omega)]}\prod_k d\omega^{(k)}.
\end{equation}
由于我们考虑的所有群都是紧致群，为了使 $U(\omega)$ 覆盖所有群元素，参数 $\omega^{(k)}$ 只需在有限区间上积分。(3.15) 中的常数 $c$ 可以用来满足归一化条件 (3.12)。剩下的需要证明的是，由 (3.14) 和 (3.15) 定义的测度确实在群内的变换下不变，即它满足 (3.11)。设 $U(\tilde\omega)$ 是由 $U(\omega)$ 与某个其他群元素 $V$ 左乘或右乘所得的群元素。参数 $\tilde\omega^{(k)}$ 与原始参数 $\omega^{(k)}$ 通过某种函数关系相联系，即我们可以写 $\omega^{(k)} = \omega^{(k)}(\tilde\omega)$。$\omega^{(k)}$ 的乘积测度按如下方式变换
\begin{equation}
  \prod_k d\omega^{(k)} = \det[J]\prod_k d\tilde\omega^{(k)},
\end{equation}
其中雅可比行列式由下式给出
\begin{equation}
  J_{kn} = \frac{\partial \omega^{(k)}}{\partial \tilde\omega^{(n)}}.
\end{equation}
$U(\tilde\omega)$ 的度规张量 $g(\tilde\omega)$ 由下式给出
\begin{equation}
\begin{split}
  g(\tilde\omega)_{nm} &= \tr\!\left[
    \frac{\partial U(\tilde\omega)}{\partial \tilde\omega^{(n)}}
    \frac{\partial U(\tilde\omega)^{\dagger}}{\partial \tilde\omega^{(m)}}
  \right] \\
  &= \tr\!\left[
    \frac{\partial U(\omega)}{\partial \omega^{(k)}}
    \frac{\partial U(\omega)^{\dagger}}{\partial \omega^{(l)}}
  \right]\frac{\partial \omega^{(k)}}{\partial \tilde\omega^{(n)}}
   \frac{\partial \omega^{(l)}}{\partial \tilde\omega^{(m)}}
  = g(\omega)_{kl}\,J_{kn} J_{lm}.
\end{split}
\end{equation}
用矩阵记号，最后一个方程写成（$T$ 表示转置）
\begin{equation}
  g(\tilde\omega) = J^T g(\omega)\,J .
\end{equation}
利用行列式的分解性质，我们从这个方程得到
\begin{equation}
  \det[g(\tilde\omega)] = \det[g(\omega)]\,\det[J]^2 .
\end{equation}
把各项综合起来，我们发现我们的群测度确实是不变的，
\begin{equation}
\begin{split}
  d\tilde U &= c\,\sqrt{\det[g(\tilde\omega)]}\prod_k d\tilde\omega^{(k)}
  = c\,\sqrt{\det[g(\omega)]}\,\det[J]\prod_k d\tilde\omega^{(k)} \\
  &= c\,\sqrt{\det[g(\omega)]}\prod_k d\omega^{(k)} = dU .
\end{split}
\end{equation}
在第一步中我们使用了 (3.20)，在第二步中使用了 (3.16)。这完成了 Haar 测度不变性的证明。

\subsection{几个 SU(3) 积分}

为了熟悉 Haar 测度的性质，我们现在讨论几个后续将用到的 SU(3) 积分。我们的讨论基于测度 $dU$ 的不变性。我们所研究的积分是基本表示中群元素 $U$ 的矩阵元 $U_{ab}$ 之积的积分。特别地，我们考虑如下积分
\begin{equation}
  \int_{SU(3)} dU\; U_{ab} = 0 ,
\end{equation}
\begin{equation}
  \int_{SU(3)} dU\; U_{ab}\, U_{cd} = 0 ,
\end{equation}
\begin{equation}
  \int_{SU(3)} dU\; U_{ab}\, (U^\dagger)_{cd} = \frac{1}{3}\,\delta_{ad}\,\delta_{bc} ,
\end{equation}
\begin{equation}
  \int_{SU(3)} dU\; U_{ab}\, U_{cd}\, U_{ef} = \frac{1}{6}\,\varepsilon_{ace}\,\varepsilon_{bdf} .
\end{equation}
分析这些积分的基本工具是关于函数 $f(U)$ 的积分的如下方程：
\begin{equation}
  \int_{SU(3)} dU\; f(U)
  = \int_{SU(3)} dU\; f(VU)
  = \int_{SU(3)} dU\; f(UW) .
\end{equation}
元素 $V$ 和 $W$ 是任意的 SU(3) 矩阵。这个关系直接来自测度的不变性 (3.11)。将 (3.26) 应用于第一个积分 (3.22) 的左端，我们得到
\begin{equation}
  \int_{SU(3)} dU\; U_{ab}
  = \int_{SU(3)} dU\; (VU)_{ab}
  = V_{ac}\int_{SU(3)} dU\; U_{cb} ,
\end{equation}
其中指标 $c$ 被求和。为了以非平凡的方式满足这个方程，必须有 $V_{ac} = \delta_{ac}$。然而，最后一个方程必须对任意的群元素 $V$ 成立，因此积分本身必须为零，从而确立了 (3.22)。用等价的方式可以证明 (3.23)。

方程 (3.24) 更有意思一些。让我们暂时在被积函数中设 $b = c$ 并对 $b$ 求和。我们得到
\begin{equation}
  \begin{split}
    \int_{SU(3)} dU\;\big[ U_{a1}(U^\dagger)_{1d} + U_{a2}(U^\dagger)_{2d} + U_{a3}(U^\dagger)_{3d} \big]
    &= \int_{SU(3)} dU\; (UU^\dagger)_{ad}\\
    &= \int_{SU(3)} dU\; \delta_{ad} = \delta_{ad} .
  \end{split}
\end{equation}
在第二步中我们使用了 $U^\dagger = U^{-1}$，在最后一步中使用了 Haar 测度的归一化 (3.12)。对于 (3.28) 的左端，重要的是认识到这三个贡献是完全等价的。例如，第一项可以通过交换第 1、2 行以及第 1、2 列而变成第二项。这一操作把 SU(3) 的一个元素变换为另一个群元素，而这两个元素在群积分中都被求和。这一等价性意味着左端的三项每一项对 (3.28) 右端的结果各贡献三分之一。于是我们确立了在 $b = c$（不求和的）情形下的公式 (3.24)。对于 $b \neq c$，我们可以再次应用 (3.26) 来证明在这种情形下积分必须为零。(3.26) 的实质是：只有当被积函数对 1 有贡献时，积分才给出非零结果，即 Haar 测度投影出被积函数对所谓的平庸表示或单重态表示的贡献。

从 (3.24) 立即得出一个以后有用的关系：
\begin{equation}
  \int dU\; \tr[VU]\,\tr[U^\dagger W] = \frac{1}{3}\,\tr[VW] .
\end{equation}
这个方程允许我们对两个基本方格乘积中出现的公共链接变量进行积分，结果是围绕两个基本方格的链接变量的迹。这一结果描绘在图 3.1 中。

\begin{figure}[htbp]
  \centering
  \includegraphics[width=0.7\textwidth]{fig31.png}
  \caption*{Fig. 3.1. \emph{对两个基本方格乘积的公共链接进行积分}.}
  \label{fig:31}
\end{figure}

本质上与 (3.24) 相同的机制也作用于积分 (3.25)：同样，被积函数的不同可能组合可以被求和以产生一个单重态。特别地，如果把被积函数的第二组指标 $b, d, f$ 设为值 $b = 1$、$d = 2$、$f = 3$，并用完全反对称的 $\varepsilon$ 张量对第一组指标 $a, c, e$ 求和，我们发现
\begin{equation}
  \varepsilon_{ace}\, U_{a1}\, U_{c2}\, U_{e3} = \det[U] = 1 ,
\end{equation}
其中我们使用了行列式的定义以及 SU(3) 的元素满足 $\det[U] = 1$ 这一事实。如果 (3.30) 的被积函数中第二组指标中有两个相等，例如 $b = d = 1$，那么用 $\varepsilon$ 张量求和对应于一个具有两列相等的行列式，因此这个和为零。如果第二组指标中的两个相对于 (3.30) 被交换，例如 $b = 2$、$d = 1$，那么这对应于两行的交换，行列式多一个负号。于是，我们得到把 (3.30) 的被积函数与 $\varepsilon$ 张量求和的如下公式
\begin{equation}
  \varepsilon_{ace}\, U_{ab}\, U_{cd}\, U_{ef} = \varepsilon_{bdf} .
\end{equation}
于是，只要第一组系数 $a, c, e$ 和第二组 $b, d, f$ 都产生非零的 $\varepsilon$ 张量贡献（$\varepsilon_{ace} \neq 0$、$\varepsilon_{bdf} \neq 0$），被积函数就对单重态有贡献。这一贡献的符号由 (3.25) 右端的 $\varepsilon$ 张量组合给出。如果这个 $\varepsilon$ 张量的乘积对给定的组合 $a, b, c, d, e, f$ 为零，那么被积函数对单重态没有贡献，积分也为零。利用与前面相同的、基于行与列同时交换的论证，人们发现所有非零贡献具有相同的权重。(3.31) 左端有六项这一事实，决定了 (3.25) 右端的因子 $1/6$。

这完成了我们对群积分的探访。积分 (3.22)、(3.23)、(3.24)、(3.25) 和 (3.29) 可以容易地推广到 SU($N$) 的情形。关于更一般的 SU($N$) 群积分的入门性论述可以在 Creutz 的著作 [1] 中找到。最后我们指出，若要比我们的第 3.4 节更详细地处理强耦合展开，特征标展开的概念是一个有力的工具。然而，特征标展开以及 SU($N$) 的一般表示理论超出了本书的范围，我们请读者参阅 [2--5]。

\section{规范不变性与规范固定}

我们已经看到，规范不变性是构造 QCD 的核心原则之一。在本节中，我们更详细地讨论选择规范的自由，并阐述规范固定在连续理论与在格点上的不同作用。我们论证，物理可观测量必须是规范不变的。我们强调，这里为纯规范理论讨论的所有规范性质都可以推广到包含费米子的完整理论。

\subsection{极大树}

让我们以一类格点所特有的规范来开始我们关于规范固定的讨论。我们证明，在某组链接变量（即所谓的极大树）中，链接变量可以被设为恒等元 1，并在规范场的积分中被略去。

我们从链接变量 $U_\mu(n)$ 的一个任意位形出发。在我们构造的开端，所有规范变换矩阵 $\Omega(n)$ 都被设为恒等元 1。我们选取一个单独的链接变量 $U_{\mu_0}(n_0)$，把链接端点处的变换矩阵 $\Omega(n_0 + \hat\mu_0)$ 设为 $U_{\mu_0}(n_0)$，保持所有其他 $\Omega(n)$ 为 1。于是我们的链接 $U_{\mu_0}(n_0)$ 被变换为
\begin{equation}
  U'_{\mu_0}(n_0)
  = \Omega(n_0)\,U_{\mu_0}(n_0)\,\Omega(n_0+\hat\mu_0)^\dagger
  = 1\,U_{\mu_0}(n_0)\,U_{\mu_0}(n_0)^\dagger = 1 .
\end{equation}
具有非平凡 $\Omega(n_0+\hat\mu_0)$ 的变换也会影响所有从格点 $n_0+\hat\mu_0$ 出发的链接变量，特别是
\begin{equation}
  U'_{\mu_1}(n_0+\hat\mu_0) = \Omega(n_0+\hat\mu_0)\,U_{\mu_1}(n_0+\hat\mu_0) .
\end{equation}
对于这些变换后的链接，我们可以重复 (3.32) 的步骤，选择它们端点处的矩阵 $\Omega$，使得这些链接也被变换为恒等元 1。整个过程可以重复进行，直到我们碰到一条链接 $U_{\mu'}(n')$，它与另一条之前已经被变换为 1 的链接相连。如果我们也想变换这条特定的链接 $U_{\mu'}(n')$，那么根据 (3.33)，我们会把另一条已经为 1 的链接从恒等元变换走。这把可以被变换为 1 的链接集合限制为一个不包含闭合回路的链接簇。这一要求把极大树定义为一组可以被变换为 1 的链接变量的极大集合（注意存在许多不同的极大树）。在图 3.2 中我们展示了一个二维子格上极大树的例子。但请注意，比极大树更小的子集也可以被规范化为 1。

\begin{figure}[htbp]
  \centering
  \includegraphics[width=0.6\textwidth]{fig32.png}
  \caption*{Fig. 3.2. \emph{二维子格上极大树的示意图。粗线表示被设为 1 并在积分中省略的链接变量}.}
  \label{fig:32}
\end{figure}

现在让我们讨论在极大树或极大树的子集上固定规范这一自由度的后果。我们考虑 (3.1)、(3.2) 和 (3.3) 中定义的某个规范不变可观测量 $O$ 的真空期望值。$O$ 的规范不变性意味着
\begin{equation}
  O[U] = O[U'] ,
\end{equation}
其中位形 $U$ 和 $U'$ 中的链接变量通过规范变换 (3.6) 相联系。我们将在下面更详细地讨论我们的可观测量的规范不变性的必要性。作用量和测度也是规范不变的，即
\begin{equation}
  S_G[U] = S_G[U'] , \qquad D[U] = D[U'] .
\end{equation}
方程 (3.34) 和 (3.35) 意味着，当把极大树中的链接设为 1 时，(3.1) 中的整个被积函数和测度保持不变。由于极大树的构造对任何特定的链接变量选择都成立，我们可以在整个积分 $D[U]$ 过程中把极大树中的链接保持为 1。由于 Haar 测度归一化为 1（见 (3.12)），极大树中的链接可以在积分中被完全省略。于是，我们可以把固定规范到极大树的过程总结如下：选取一棵极大树，或者一个不含闭合回路的较小的链接集合，并把该集合中的链接设为 1。随后，对所有不包含在所选集合作中的链接变量进行积分。无论我们是否固定规范，规范不变可观测量的期望值都不受影响。

规范固定在格点上的作用与它在连续理论中扮演的角色非常不同。在格点上固定规范是我们为实现某些计算简化或使可观测量的诠释更加透明而可以采取的一个步骤（见下面关于威尔逊圈的讨论）。然而，固定规范并不是使算符的真空期望值得到良好定义并可计算的一个必要步骤。在连续理论中情形不同：在那里，胶子作用量二次部分的核（即逆胶子传播子）具有由纯规范引起的零模，也就是对平凡规范场位形的规范变换。这些模导致胶子传播子中出现奇异性，必须通过选取特定的规范来消除。只有在那一步之后，连续理论才能在微扰论中得到良好定义。

我们已经提到，通过极大树固定规范的程序是格点表述所特有的。然而，存在某种规范，即所谓的时序规范，它在连续理论中有相应的对应物。在时间方向为无穷大的格点上，我们可以把所有类时链接设为 1，即我们可以设
\begin{equation}
  U_4(n) = 1 \quad \forall n .
\end{equation}
我们指出，这不是一棵极大树，而是一个更小的集合。在（欧几里得）连续理论中，时序规范由 $A_4(x) = 0$ 定义，由于关系式 $U_4 = \exp(\im a A_4)$，这与格点定义相符。时序规范对于连续规范理论的哈密顿表述很重要（参见任何量子场论教材，例如 [6--8]）。当我们在引入和讨论威尔逊圈时，我们将利用时序规范。

\subsection{其他规范}

在连续理论中使用的其他规范在格点上也起着作用，特别是当人们把格点结果与在特定规范（例如朗道规范、库仑规范）中完成的连续计算进行匹配时。在对禁闭机制的格点研究中，也会使用特殊的规范，例如最大阿贝尔规范。这些规范不是通过极大树实现的，而是通常通过使某个泛函取极值来实现。作为说明这类规范如何构造的一个例子，我们简要讨论朗道规范，在讨论重整化的第 11 章中我们将需要它。

在连续理论中，朗道规范由条件
\begin{equation}
  \partial_\mu A_\mu(x) = 0 .
\end{equation}
定义。关键的一步是把这一在连续理论中以代数取值规范场 $A_\mu(x)$ 表述的规则，转译为一个可以在格点上使用的规则，在格点上我们处理的是群取值链接变量 $U_\mu(n)$。第一步，我们证明要求朗道规范条件等价于求泛函
\begin{equation}
  W[A] = \sum_{\mu=1}^{4}\int \dd^4 x\, \tr\!\left[A_\mu(x)^2\right] ,
\end{equation}
关于规范变换的极值。如果 $W[A]$ 处于极值，那么当我们施加由 $\Omega(x) = \exp(\im\,\varepsilon\, H(x))$ 定义的无穷小规范变换时，它必须保持不变。注意，$H(x)$ 是一个无迹厄米矩阵，因此 $\Omega(x)^\dagger = \exp(-\im\,\varepsilon\, H(x))$。一行代数运算表明，对于我们的无穷小变换，规范变换 (2.11) 化为
\begin{equation}
  A_\mu(x) \to A_\mu(x) + \varepsilon\,\big(\im\,[H(x), A_\mu(x)] - \partial_\mu H(x)\big) + \order{\varepsilon^2} .
\end{equation}
把变换后的规范场代入泛函 $W[A]$，人们发现它按如下方式变化
\begin{equation}
\begin{split}
  W[A] &\to W[A] - 2\varepsilon\sum_{\mu=1}^{4}\int \dd^4 x\,
    \tr\!\left[A_\mu(x)\,\partial_\mu H(x)\right] + \order{\varepsilon^2} \\
  &= W[A] + 2\varepsilon\sum_{\mu=1}^{4}\int \dd^4 x\,
    \tr\!\left[(\partial_\mu A_\mu(x))\, H(x)\right] + \order{\varepsilon^2} ,
\end{split}
\end{equation}
其中迹的循环性使得 (3.39) 中对易子的贡献为零。在第二步中我们进行了分部积分（假设 $H(x)$ 在大 $x$ 处为零）。如果 $W[A]$ 处于极值，那么 (3.40) 中 $\order{\varepsilon}$ 的项必须为零。由于 $H(x)$ 可以选为任意的无迹厄米矩阵，这一陈述等价于 (3.37)。

把 $W[A]$ 转译为相应的格点规则很简单。我们可以考虑泛函
\begin{equation}
  W_{\mathrm{lat}}[U] = -a^2\sum_{n\in\Lambda}\sum_{\mu=1}^{4}
    \tr\!\left[U_\mu(n) + U_\mu(n)^\dagger\right] ,
\end{equation}
把 $U_\mu(n) = \exp(\im a A_\mu(n))$ 对小 $a$ 展开后，它就变成 $W[A]$（附加的常数是无关紧要的）。

目标是找到一个规范变换 $\Omega$，对于给定的规范位形，它把 $W_{\mathrm{lat}}[U]$ 移动到极小值。更明确地说，我们必须求
\begin{equation}
  F[\Omega] = -a^2\sum_{n\in\Lambda}\sum_{\mu=1}^{4}
    \tr\!\left[\Omega(n)U_\mu(n)\Omega(n+\hat\mu)^\dagger
    + \Omega(n+\hat\mu)\,U_\mu(n)^\dagger\,\Omega(n)^\dagger\right] ,
\end{equation}
作为固定 $U_\mu(n)$ 时矩阵 $\Omega(n)$ 的函数。这是一个大规模优化问题，可以用诸如超松弛或模拟退火之类的方法来求解（参见，例如，[9]）。一个特别的问题在于矩阵 $\Omega(n)$ 必须属于 SU(3) 这一约束。我们将在第 4 章讨论如何在数值上生成这样的矩阵。我们指出，对这里为朗道规范提出的方法作一个直接的修改（仅对 $\mu$ 从 1 到 3 求和），就可以在格点上实现库仑规范条件 $\sum_{j=1}^{3}\partial_j A_j = 0$。关于格点上规范固定程序的综述，参见 [10]。

\subsection{可观测量的规范不变性}

在我们关于极大树的讨论中，我们已经假设物理可观测量 $O$ 是规范不变的，即它们满足 (3.34)。我们现在更详细地讨论这一要求。

通过检视积分 (3.22)，人们看到，试图把单个链接变量用作可观测量的做法是失败的，因为它的期望值为零。这样一个单独的链接变量是非规范不变泛函 $O[U]$ 的一个例子，即对于由规范变换相联系的位形 $U$ 和 $U'$，它不满足 $O[U] = O[U']$。任何非规范不变泛函的期望值都等于它在规范群上的平均值的期望。为看清这一点，我们施加规范变换 $U_\mu(n) \to \Omega(n)U_\mu(n)\Omega(n+\hat\mu)^\dagger$，并得到任意可观测量的期望值
\begin{equation}
  \begin{split}
    \ev{O} &= \frac{1}{Z}\int D[U]\,e^{-S_G[U]}\,O[U]\\
    &= \frac{1}{Z}\int D[\Omega U\Omega^\dagger]\,e^{-S_G[\Omega U\Omega^\dagger]}
      \,O[\Omega U\Omega^\dagger]\\
    &= \frac{1}{Z}\int D[U]\,e^{-S_G[U]}\,O[\Omega U\Omega^\dagger] ,
  \end{split}
\end{equation}
其中在最后一步我们同时使用了作用量的规范不变性和 Haar 测度的不变性。由于值 $\Omega(n)$ 可以任意选取，我们甚至可以在规范群上取平均：
\begin{equation}
  \ev{O} = \frac{1}{Z}\int D[U]\,e^{-S_G[U]}
    \int D[\Omega]\,O[\Omega U\Omega^\dagger]
  = \ev{\int D[\Omega]\,O[\Omega U\Omega^\dagger]} .
\end{equation}
循着这些论证的思路，Elitzur [11] 证明了局域规范对称性不可能自发破缺。

根据 (3.44)，对我们的单个链接变量的例子，我们发现
\begin{equation}
  \ev{U_\mu(n)}
  = \int d\Omega(n)\,\ev{\Omega(n)U_\mu(n)}
  = \int d\Omega(n)\,\Omega(n)\,\ev{U_\mu(n)} = 0 .
\end{equation}
我们指出，对某些计算而言，使用非规范不变的可观测量并在规范固定之后再对其进行求值，确实是有利的。如所提及的那样，当格点结果要与在特定规范中完成的连续计算进行比较时，这一程序是有用的。

\section{威尔逊圈与波利亚科夫圈}

在本节中，我们介绍允许人们确定两个静态色源之间势的可观测量。这些可观测量就是所谓的威尔逊圈和波利亚科夫圈，我们首先引入它们，稍后再给出它们的诠释。

\subsection{威尔逊圈的定义}

在上一节中我们讨论了物理可观测量必须是规范不变的。一个仅由规范场构成、且规范不变的原型对象，是我们在 (2.46) 中引入的沿闭合回路的链接变量乘积的迹，
\begin{equation}
  L[U] = \tr\!\left[\prod_{(n,\mu)\in L} U_\mu(n)\right] .
\end{equation}
这里 $L$ 是格点上的一个闭合链接回路，(3.46) 中的乘积遍历所有这些链接。我们现在引入的威尔逊圈就属于这种类型。一个威尔逊圈 $W_L$ 由四个部分构成，即两条所谓的威尔逊线 $S(m, n, n_t)$、$S(m, n, 0)$ 和两个时间方向输运子 $T(n, n_t)$、$T(m, n_t)$。威尔逊线 $S(m, n, n_t)$ 沿某条路径 $C_{m,n}$ 连接两个空间点 $m$ 和 $n$，其中所有链接变量的时间分量都固定为 $n_t$，
\begin{equation}
  S(m, n, n_t) = \prod_{(k,j)\in C_{m,n}} U_j(k, n_t) .
\end{equation}
时间方向输运子 $T(n, n_t)$ 是时间方向上由 $n_t$ 个链接变量组成的直线，所有这些变量都位于空间位置 $n$ 处，
\begin{equation}
  T(n, n_t) = \prod_{j=0}^{n_t-1} U_4(n, j) .
\end{equation}
把四个部分彼此相连就给出一个闭合回路 $L$，
\begin{equation}
  L: \quad (m, n_t) \xrightarrow{S} (n, n_t) \xrightarrow{T^\dagger} (n, 0)
     \xrightarrow{S^\dagger} (m, 0) \xrightarrow{T} (m, n_t) .
\end{equation}
取迹就得到威尔逊圈 $W_L$，
\begin{equation}
  W_L[U] = \tr\!\left[S(m, n, n_t)\,T(n, n_t)^\dagger\,S(m, n, 0)^\dagger\,
    T(m, n_t)\right]
  = \tr\!\left[\prod_{(k,\mu)\in L} U_\mu(k)\right] .
\end{equation}
如果 $S(m, n, n_t)$ 中使用的回路段 $C_{m,n}$ 是一条直线，我们就称之为平面威尔逊圈。注意，只有当 $m$ 和 $n$ 落在同一条坐标轴上时才可能出现这种情况。否则威尔逊圈称为非平面威尔逊圈。图 3.3 展示了一个平面圈和一个非平面圈的示例。

\begin{figure}[htbp]
  \centering
  \includegraphics[width=0.8\textwidth]{fig33.png}
  \caption*{Fig. 3.3. \emph{平面（左图）与非平面（右图）威尔逊圈的示例。水平方向为时间}.}
  \label{fig:33}
\end{figure}

\subsection{时序规范}

在我们能够讨论威尔逊圈的物理诠释之前，我们必须讨论规范理论的一个特殊之处。如果想对规范场作用量 (2.17) 求正则动量 (1.45)，人们会发现，对于时间分量 $A_4$，正则动量为零。原因是场强张量 $F_{\mu\nu}(x) = -\im\,[D_\mu(x), D_\nu(x)]$ 不包含 $A_4$ 对时间的导数。摆脱这一困境的一个可能办法是使用满足如下条件的规范
\begin{equation}
  A_4(x) = 0 ,
\end{equation}
即前面提到的时序规范。我们指出，仅仅陈述 (3.51) 远不能公正对待规范理论量子化所涉及的微妙之处，关于这一问题我们请读者参阅诸如 [6--8] 之类的场论著作。在格点上，时序规范对应于条件 (3.36)。

我们强调，在下面的讨论中，我们使用时序规范只是为了找到威尔逊圈的物理诠释。对于期望值的实际计算，我们并不需要固定规范。无论我们是否固定规范，威尔逊圈期望值的结果当然是相同的。

\subsection{威尔逊圈的物理诠释}

在上一段讨论的时序规范 (3.36) 中，时间方向输运子变为平庸的，
\begin{equation}
  T(n, n_t) = \prod_{j=0}^{n_t-1} U_4(n, j) = 1 ,
\end{equation}
并且我们得到如下恒等式链
\begin{equation}
  \ev{W_L} = \ev{W_L}_{\mathrm{temp}}
  = \ev{\tr\!\left[S(m, n, n_t)\,S(m, n, 0)^\dagger\right]}_{\mathrm{temp}} ,
\end{equation}
其中在第一步中我们使用了以下事实：固定规范时，规范不变可观测量的期望值保持不变。在第二步中我们使用了 $W_L[U] = \tr[S(m, n, n_t)\,S(m, n, 0)^\dagger]$，这来自 (3.50) 中威尔逊圈的定义以及时序规范 (3.52) 中时间方向输运子的平庸性。

(3.53) 中使用的时序规范清楚地表明，威尔逊圈是位于时间片 $n_t$ 和 0 的两条威尔逊线 $S(m, n, n_t)$ 与 $S(m, n, 0)$ 的关联函数。因此，我们可以用第一条关键方程的形式 (1.21) 来诠释这个关联函数。据此，当我们的欧几里得格点的总时间范围 $T$ 很大时，关联函数表现为（对 $a$、$b$ 求和）
\begin{equation}
  \ev{\tr\!\left[S(m, n, n_t)\,S(m, n, 0)^\dagger\right]}_{\mathrm{temp}}
  = \sum_k \bra{0}\tilde S(m,n)_{ab}\ket{k}\,
    \bra{k}\tilde S(m,n)^\dagger_{ba}\ket{0}\, e^{-tE_k} ,
\end{equation}
其中欧氏时间参数 $t$ 通过 $t = a\,n_t$ 与 $n_t$ 相联系，$a$ 是格点间距。(3.54) 中的求和遍历所有与 $\tilde S(m,n)^\dagger\ket{0}$ 有非零交叠的态 $\ket{k}$。

在下一段中我们将论证，具有非零交叠的态 $\ket{k}$ 是描述位于空间位置 $m$ 和 $n$ 的静态夸克-反夸克对的态。因此，在 (3.54) 中，能量最低 $E_1$ 的项预期就是描述我们静态夸克-反夸克对的态。更高的态可以是，例如，这个对加上具有真空量子数的额外粒子-反粒子组合。于是能量 $E_1$ 被等同于夸克-反夸克对的能量，即夸克空间分离为 $r$ 时的静态势 $V(r)$，
\begin{equation}
  E_1 = V(r) \qquad \text{with} \quad r = a\,|m - n| .
\end{equation}
把 (3.53)、(3.54) 和 (3.55) 结合起来，我们得到
\begin{equation}
  \ev{W_L} \propto e^{-tV(r)}\left(1 + \order{e^{-t\Delta E}}\right)
  = e^{-n_t a V(r)}\left(1 + \order{e^{-n_t a \Delta E}}\right) .
\end{equation}
于是我们发现，我们可以从威尔逊圈的大 $n_t$ 行为来计算静态夸克-反夸克势。(3.56) 中的修正项被指数压低，其中 $\Delta E$ 是 $V(r)$ 与夸克-反夸克对第一激发能级之差。

我们强调，威尔逊圈是有取向的。然而，反转取向在代数层面上就是复共轭（比较 (2.34)），它仅仅对应于夸克与反夸克的交换。因此，两种取向对于确定势 $V(r)$ 具有同样的作用。

我们引入的威尔逊圈不一定是平面的。在图 3.3 中我们展示了两个威尔逊圈，一个是平面的（左图），一个是非平面的（右图）。两个圈都有 $n_t = 5$（水平方向为时间）。平面圈的 $r = 3a$，非平面圈的 $r = \sqrt{3^2 + 1}\,a = \sqrt{10}\,a$。因此，利用非平面威尔逊圈，我们不仅可以在 $r$ 为 $a$ 的整数倍的距离处计算势 $V(r)$，还可以在中间点计算。非平面威尔逊圈还允许人们研究在趋近连续极限时旋转不变性是否最终得以恢复。

威尔逊圈也可以用作纯胶子束缚态（即所谓的胶球）的算符。相应欧几里得关联函数的指数衰减使得人们能够提取这些态的质量谱（参见第 6 章以及 [12] 中的讨论）。

\subsection{威尔逊线与夸克-反夸克对}

为了完成我们对威尔逊圈的物理诠释，我们还需要证明态 $\tilde S(m,n)^\dagger_{ba}\ket{0}$ 确实与夸克-反夸克对有交叠。然而，关于 $S(m, n)$ 是在大夸克质量极限下描述夸克-反夸克对的正确表达式的真正推导，只有到第 5 章在给出格点上夸克场的最终定义之后才会给出。在那里我们将证明，在无穷重夸克极限下夸克传播子化为 $S(m, n)$。目前我们满足于证明 $S(m, n)$ 在规范变换下具有与夸克-反夸克对相同的变换性质。

根据我们在第 2.1 节的讨论，位于空间位置 $am$、$an$ 的夸克-反夸克对由场的乘积描述
\begin{equation}
  Q(m, n)_{ab} \equiv \psi(m)^{\alpha}_{a}\,\psi(n)^{\beta}_{b} .
\end{equation}
夸克场携带旋量（$\alpha$、$\beta$）和色（$a$、$b$）指标。然而，这里我们不关心势对旋量指标的依赖，在 $Q(m, n)_{ab}$ 的定义中我们忽略它们。$Q(m, n)_{ab}$ 不是规范不变的。根据 (2.7)，它在规范变换下按如下方式变换
\begin{equation}
  Q(m, n)_{ab} \to \Omega(m)_{aa'}\,Q(m, n)_{a'b'}\,\Omega(n)^\dagger_{b'b} .
\end{equation}
从链接变量规范变换性质 (2.33) 之后的讨论中我们知道，(2.43) 中定义的链接变量乘积 $P[U]$ 恰按所要求的方式变换（比较 (2.44)），
\begin{equation}
  S(m, n)_{ab} \to \Omega(m)_{aa'}\,S(m, n)_{a'b'}\,\Omega(n)^\dagger_{b'b} .
\end{equation}
于是我们至少验证了威尔逊线具有与夸克-反夸克对相同的变换性质。

\subsection{波利亚科夫圈}

让我们以讨论威尔逊圈的一种变形——即所谓的波利亚科夫圈 [13]（也称为热威尔逊线）——来结束本节。这里我们采用规范场在时间方向上周期的边界条件。我们把威尔逊圈的时间范围 $n_t$ 在我们格点上取为尽可能大，即设 $n_t = N_T$，其中 $N_T$ 是时间方向上格点的总数。于是威尔逊圈的空间部分（图 3.3 中的非水平线）相互叠置，但取向相反。由于周期边界条件，我们不能把所有时间链接都规范变换为 1。然而，我们可以把圈的空间部分规范化为 1。于是威尔逊圈约化为时间链接变量的两条不连通路径（比较 (3.48)）$T(m, N_T)$、$T(n, N_T)^\dagger$，它们位于空间中的两个位置 $m$ 和 $n$ 处。这两条路径都缠绕着格点的时间方向，但取向相反。

我们可以分别对这两条回路取迹，使这个新的可观测量变为规范不变的。这不过是色指标的一种重新排列，并不改变可观测量的诠释。这样我们就引入了所谓的波利亚科夫圈
\begin{equation}
  P(m) = \tr\!\left[\prod_{j=0}^{N_T-1} U_4(m, j)\right] ,
\end{equation}
它作为闭合回路上的迹是规范不变的。我们现在可以放弃我们的特殊规范，并得到（$r = a|m - n|$）
\begin{equation}
  \ev{P(m)P(n)^\dagger} \propto e^{-N_T a V(r)}\left(1 + \order{e^{-N_T a\Delta E}}\right) .
\end{equation}
我们在第 4 章给出的静态势的数值计算基于波利亚科夫圈，即基于两个方程 (3.60) 和 (3.61)。分别通过威尔逊圈或波利亚科夫圈所定义的势的同一性尚未被严格证明；已有一个证明表明，从波利亚科夫圈关联函数导出的弦张力以威尔逊圈的弦张力为上界 [14]。

引入波利亚科夫圈的另一种动机方式是把流 $j_\mu$ 耦合到场 $A_\mu$，即定义算符 $O = \tr\,P\exp\!\big(\im\int \dd^4 z\, j_\mu(z)\, A_\mu(z)\big)$。利用对应于位置 $x$ 处静态电荷的流，即 $j_\mu(z) = (0, 0, 0, 1)\,\delta(z - x)$（欧几里得度规），在格点上算符 $O$ 就转译为我们的波利亚科夫圈。

最后我们指出，单个波利亚科夫圈的真空期望值 $\ev{P(n)}$ 也是一个重要的变量。正如我们将在第 12 章讨论的那样，它是有限温度下胶动力学退禁闭相变的序参量。

\section{静态夸克势}

在把威尔逊圈引入为静态夸克势 $V(r)$ 的可观测量之后，我们现在讨论 $V(r)$ 的一般形式。作为第一步，我们在强耦合 $g$ 的极限下（对应于小 $\beta$——比较 (3.5)）计算势，并发现这产生一个线性上升的项。随后，我们论证，对于小的耦合 $g$，我们得到从电动力学熟知的 $1/r$ 势。于是，我们发现静态 QCD 势可以用下式参数化
\begin{equation}
  V(r) = A + \frac{B}{r} + \sigma r .
\end{equation}
由于夸克之间的力是 $V(r)$ 的导数，常数 $A$ 只是能量的一个无关紧要的归一化。(3.62) 中的第二项是强度为 $B$ 的势的库仑部分。最后，第三项贡献是一个线性上升的项，实常数 $\sigma$ 就是所谓的弦张力。根据 QCD 唯象学，人们预期 $\sigma \approx 900\ \mathrm{MeV}/\mathrm{fm}$。在为 (3.62) 中的各项提供了证据之后，我们讨论静态 QCD 势的物理意义。

\subsection{威尔逊圈的强耦合展开}

为了证明 (3.62) 中线性上升项的存在，我们在强耦合极限（即大 $g$、小 $\beta$）下计算威尔逊圈的真空期望值。更明确地说，我们计算
\begin{equation}
  \ev{W_C} = \frac{1}{Z}\int D[U]\,
    \exp\!\left[-\frac{\beta}{3}\sum_P \Re\,\tr[1 - U_P]\right]
    \tr\!\left[\prod_{l\in C} U_l\right] .
\end{equation}
对于这一计算，我们使用一个简化的记号：求和遍历所有基本方格 $P$，其中每个基本方格只按两种可能取向之一来计数。对 $l$ 的乘积遍历定义威尔逊圈的回路 $C$ 中所包含的所有链接变量。这个表达式可以改写为
\begin{equation}
\begin{split}
  \ev{W_C} &= \frac{1}{Z'}\int D[U]\,
    \exp\!\left[\frac{\beta}{3}\sum_P \Re\,\tr[U_P]\right]
    \tr\!\left[\prod_{l\in C} U_l\right] \\
  &= \frac{1}{Z'}\int D[U]\,
    \exp\!\left[\frac{\beta}{6}\sum_P \left(\tr[U_P] + \tr[U_P^\dagger]\right)\right]
    \tr\!\left[\prod_{l\in C} U_l\right] .
\end{split}
\end{equation}
在第一步中，我们把常数因子 $\exp\!\left(-\frac{\beta}{3}\sum_P \Re\,\tr[1]\right)$ 从玻尔兹曼因子 $\exp(-S)$ 中分离出来。完全相同的常数因子也出现在配分函数 $Z$ 中，我们把分子和分母中的两个因子相互抵消。不含这个因子的配分函数记作 $Z'$。在第二步中我们使用
\begin{equation}
  \Re\,\tr[U_P] = \frac{1}{2}\left(\tr[U_P] + \tr[U_P^\dagger]\right) .
\end{equation}
我们强调，根据 (2.48)，基本方格变量 $U_P$ 的厄米共轭等价于反转基本方格的取向。因此，在 (3.64) 的第二行中我们明确写出了基本方格变量 $U_P$ 的两种取向，这导致一个额外的因子 $1/2$。

在 (3.64) 的形式下，我们现在可以讨论强耦合（小 $\beta$）下威尔逊圈期望值的展开。特别地，我们利用指数函数的泰勒展开把 (3.64) 的玻尔兹曼因子按 $\beta$ 展开，
\begin{equation}
  \begin{split}
    \exp\!\left[\frac{\beta}{6}\sum_P \left(\tr[U_P] + \tr[U_P^\dagger]\right)\right]
    &= \sum_{i,j=0}^{\infty}\frac{1}{i!\,j!}\left(\frac{\beta}{6}\right)^{i+j}\\
    &\qquad \left(\sum_P \tr[U_P]\right)^{i}
      \left(\sum_P \tr[U_P^\dagger]\right)^{j} .
  \end{split}
\end{equation}
注意，在这个展开中我们把顺时针取向的基本方格 $U_P^\dagger$ 和逆时针取向的基本方格 $U_P$ 的贡献分离开来。这一点很重要，因为对于展开中的领头项，只有那些取向与威尔逊圈相反的基本方格才起贡献。

对于归一化因子 $Z'$，确定小 $\beta$ 展开中的领头贡献是直截了当的。(3.66) 中 $i = j = 0$ 的第一项就已经对积分给出非零贡献，我们得到（利用 Haar 测度的归一化）
\begin{equation}
  \begin{split}
    Z' &= \int D[U]\,\exp\!\left[\frac{\beta}{6}\sum_P
      \left(\tr[U_P] + \tr[U_P^\dagger]\right)\right]\\
    &= \int D[U]\,\left(1 + \order{\beta}\right)
    = 1 + \order{\beta^2} .
  \end{split}
\end{equation}
(3.64) 的分子展开不那么直截了当。如果只保留玻尔兹曼因子展开中的领头项，那么构成可观测量的链接变量乘积 $\prod_{l\in C} U_l$ 会产生 (3.22) 类型的积分，它们全部为零。因此，为了找到展开中首个非零项，我们必须把玻尔兹曼因子按小 $\beta$ 展开。这从指数中带下额外的链接变量，这样我们就可以饱和对链接的积分，从而得到 (3.24) 类型的非零贡献。在次领头项中，(3.25) 也有贡献。

如果我们把威尔逊圈的回路 $C$ 视为一个 $n_r \times n_t$ 的链接矩形，那么这个回路所张的最小面积 $A_C$ 包含 $n_A = n_r n_t$ 个基本方格（注意 $n_r$、$n_t$、$n_A$ 都是正整数）。物理面积 $A_C$ 与威尔逊圈在物理单位 $a n_r$、$a n_t$ 下的延伸通过 $A_C = a^2 n_A = a n_r\, a n_t$ 相联系。回顾 (3.29) 和图 3.1，我们发现，只有当圈中每个链接变量 $U_\mu(n)$ 与其共轭伙伴 $U_\mu(n)^\dagger$ 配对时，才出现非零贡献。由于我们的作用量中含有基本方格，这个过程必须持续进行，直到我们用从玻尔兹曼因子展开得到的 $n_A$ 个基本方格填满回路 $C$。我们在图 3.4 中描绘了这一贡献。

注意，用于填充回路的那些基本方格必须具有与威尔逊圈相反的取向。只有这样，每个链接处的贡献才具有 (3.24) 中被积函数的形式。由于我们至少需要指数中的 $n_A = n_r n_t$ 个基本方格，指数展开 (3.66) 中所需要的项是 $n_A$ 阶的。这个领头项显式地写为（注意，(3.66) 中两种取向里只有与威尔逊圈相反的那一种起贡献）
\begin{equation}
\begin{split}
  &\int D[U]\;\frac{1}{n_A!}\left(\frac{\beta}{6}\right)^{n_A}
    \left(\sum_P \tr[U_P^\dagger]\right)^{n_A}
    \tr\!\left[\prod_{l\in C} U_l\right] \\
  &\qquad = \left(\frac{\beta}{6}\right)^{n_A}\int D[U]\;
    \left(\sum_{P\in A_C} \tr[U_P^\dagger]\right)
    \tr\!\left[\prod_{l\in C} U_l\right] \\
  &\qquad = \left(\frac{\beta}{6}\right)^{n_A}\tr[\mathbb{1}]
    \left(\frac{1}{3}\right)^{n_A}
  = 3\,\exp\!\left[n_A \ln\!\left(\frac{\beta}{18}\right)\right] .
\end{split}
\end{equation}
在 (3.68) 的第一步中，我们对具有正确取向（$U_P^\dagger$）的基本方格变量之和展开 $n_A$ 次幂。我们只保留最小面积 $A_C$ 内每个基本方格 $P$ 都被匹配的 $U_P^\dagger$ 所占据的那些项。这样的乘积恰好有 $n_A!$ 个，因此因子 $1/n_A!$ 被抵消。$n_A$ 次幂展开中的所有其他项都为零，因为它们产生 (3.22) 类型的积分。在第二步中，我们利用图 3.1 中所描绘形式的积分 (3.29) 来求值非零项。我们先把威尔逊圈内部的基本方格变量粘合成行，再把这些行粘合成图 3.4 中的整个内部方块。另一个规范积分把这个方块与威尔逊圈取向相反的外回路联系起来。人们发现这些步骤给出 (3.68) 中的因子 $(1/3)^{n_A}$。把 (3.67) 和 (3.68) 结合起来，我们得到
\begin{equation}
  \ev{W_C} = 3\,\exp\!\left[n_A \ln\!\left(\frac{\beta}{18}\right)\right]
    \left(1 + \order{\beta}\right)
  = 3\,\exp\!\left[n_r n_t \ln\!\left(\frac{\beta}{18}\right)\right]
    \left(1 + \order{\beta}\right) .
\end{equation}
根据 (3.56)，这个表达式必须与渐近形式相比较，即对于大的 $t = a n_t$，我们有
\begin{equation}
  \ev{W_C} \propto \exp\left(-a n_t V(r)\right) .
\end{equation}
于是我们得出结论：在强耦合极限下（注意 $r = a n_r$）

\begin{figure}[htbp]
  \centering
  \includegraphics[width=0.8\textwidth]{fig34.png}
  \caption*{Fig. 3.4. \emph{威尔逊圈强耦合（小 $\beta$）展开中的领头贡献。外部逆时针取向的矩形是威尔逊圈，较小的顺时针取向的方块是来自作用量的单个基本方格项}.}
  \label{fig:34}
\end{figure}
\begin{equation}
  V(r) = \sigma r ,
\end{equation}
其中弦张力 $\sigma$ 由领头阶表达式给出
\begin{equation}
  \sigma = -\frac{1}{a^2}\ln\!\left(\frac{\beta}{18}\right)\left(1 + \order{\beta}\right) .
\end{equation}
我们指出，产生弦张力的更高阶修正相对容易。然而，由于强耦合展开在现代格点规范理论中并不起核心作用（我们很快会看到，我们实际关心的是趋近弱耦合），我们将不讨论更高阶项的计算以及这个级数的收敛性质。\footnote{Wilson 的原始推导 [15] 已在 [16] 中被严格证明。关于初等的表述我们请读者参阅 Creutz 的著作 [1]，关于更高级的计算请参阅 [17]。} 相反，我们再次强调，我们用相对简单的展开就从格点表述中提取出了线性上升的势。下面我们将讨论，势中的这一项如何给出禁闭这一重要特征。然而，在我们讨论物理意义之前，让我们先给出静态 QCD 势参数化 (3.62) 中存在库仑型项的论证。

\subsection{静态夸克势的库仑部分}

库仑型相互作用的存在可以从连续理论中小耦合常数 $g$ 时胶子作用量的行为直接看出。这个论证最好在 (2.18) 和 (2.19) 中引入的规范场归一化下给出，那里规范场 $A_\mu$ 被重新标度了因子 $1/g$。在这种形式下，欧几里得连续规范场作用量写成对色分量求和的形式（比较 (2.23) 和 (2.24)）
\begin{equation}
  S_G[A] = \frac{1}{4}\sum_{i=1}^{8}\int \dd^4 x\,
    F^{(i)}_{\mu\nu}(x)\,F^{(i)}_{\mu\nu}(x) ,
\end{equation}
其中单个色分量的场强张量由下式给出
\begin{equation}
  F^{(i)}_{\mu\nu}(x)
  = \partial_\mu A^{(i)}_{\nu}(x) - \partial_\nu A^{(i)}_{\mu}(x)
  - g\, f^{ijk} A^{(j)}_{\mu}(x)\,A^{(k)}_{\nu}(x) .
\end{equation}
从 (3.74) 可以清楚地看到，QCD 特有的自相互作用项都乘以耦合常数 $g$。因此，当把这个耦合常数趋于零时，场强张量 (3.74) 就约化为它的阿贝尔对应物，即它具有从 QED 熟知的形式。于是，作用量 (3.73) 变成对每个色分量的 QED 型相互作用之和。由于我们知道在 QED 中静态势是库仑型的，\footnote{关于这一事实在路径积分形式体系中的初等推导，参见 [18]。} 我们得出结论，$1/r$ 项应当存在于静态 QCD 势的参数化 (3.62) 中。

\subsection{静态 QCD 势的物理意义}

让我们简要勾勒静态 QCD 势的物理意义，特别是线性项的作用。静态夸克-反夸克对之间势中的线性上升项意味着，当人们试图把两个组分拉开时，能量持续线性上升。因此，夸克和反夸克被禁闭在一个强束缚的介子态中。类似地，正如我们后面将展示的，三个夸克的组合也被束缚，形成一个重子。只有像强子这样的色中性组合才是可观测对象，这一现象是禁闭的一种可能定义。

导致线性上升项的物理机制是两个源之间通量管的形成。在 QED 中，由于规范场没有自相互作用，源与汇之间的场线在空间中铺展开来。在 QCD 中，胶子的强自相互作用阻止了这种行为，场被挤压进一根狭窄的管或弦中，从而产生线性上升。线性上升势的直接实验证据可以在把强子质量作为其总自旋的函数作图并发现线性行为时看到（参见，例如，[19] 中的讨论）。由于对于线性上升的势，能量随角动量线性上升，这一实验发现证实了 (3.62) 中的线性项。

到目前为止，我们的讨论仅基于从纯胶动力学得到的静态势。当然，尚未纳入我们讨论的夸克也起着重要作用。在具有动力学夸克的完整理论中，粒子-反粒子产生与湮灭的过程变得重要。特别是，如果夸克和反夸克被拉到足够远，能量就变得足够大，足以产生一个夸克-反夸克对，它可能与最初的两个组分重新组合而形成两个介子。这一现象称为弦断裂，可以在格点上进行研究。

总而言之，我们发现格点 QCD 对禁闭非常友好。这一性质在强耦合极限下很容易证明，但正如我们将在第 4 章看到的，在较弱耦合的数值计算中提取 QCD 势也相对简单。

\section{利用静态势设定标度}

在格点形式体系中，作用量以 $\hbar$ 为单位给出，所有可观测量都是无量纲的。只有把它们与物理量联系起来，我们才能引入标度参数。在格点上计算的这样一种无量纲量的例子是格点间距 $a$ 与某个质量 $M$ 的乘积 $aM$，它可以从像 (1.21) 那样的关联函数的指数衰减来确定。把 $M$ 等同于一个物理质量，就使我们能够以物理单位确定格点常数 $a$。

在本节中，我们讨论另一种设定标度的方法，它基于静态夸克势 $V(r)$。这使我们能够把格点间距 $a$ 与逆规范耦合 $\beta$ 联系起来。我们定义某个物理距离 $r_0$，即所谓的索末菲参数 [20]，它是与静态势相关的一个特征长度标度。以物理单位计，这个距离是 $r_0 \approx 0.5$ 费米（1 费米 $= 1$ fm $= 10^{-15}$ m）。这个距离使我们能够简单地通过数 $r = 0$ 与 $r = r_0$ 之间的格点数，从 $V(r)$ 确定格点间距 $a$。在本节结束时，我们将能够讨论纯格点规范理论的真正连续极限。

\subsection{静态势数值数据的讨论}

让我们以讨论静态势的数值数据来开始我们的介绍。此刻我们假设数据已经被计算出来，因为我们只在第 4 章才引入这种计算的实际技术。我们假设已经计算了大小为 $r \times t$ 的平面威尔逊圈的期望值 $\ev{W_C}$，即回路 $C$ 是一个 $r \times t$ 矩形。空间距离 $r$ 和欧氏时间 $t$ 通过格点间距 $a$ 与整数 $n$ 和 $n_t$ 相联系，
\begin{equation}
  r = n a , \qquad t = n_t a .
\end{equation}
根据 (3.56)，威尔逊圈的真空期望值通过下式与静态势相联系（这里我们省略了 (3.56) 中明确显示的指数压低修正）
\begin{equation}
  \ev{W_C} = C\,\exp\left(-t V(r)\right)
  = C\,\exp\left(-n_t a\,V(na)\right) .
\end{equation}
由于在这个阶段格点间距 $a$ 仍然未知，我们能从 $\ev{W_C}$ 的数值数据中提取的只是不同 $n$ 处乘积 $aV(an)$ 的值。实际上，这是通过根据 (3.76) 对不同 $n_t$（但固定 $n$）的数据进行双参数拟合来完成的。拟合参数是 $C$ 和 $aV(an)$。这个过程可以对不同的 $n$ 值重复进行，结果就是一组作为 $n$ 之函数的 $aV(an)$ 数值数据。

在图 3.5 中，我们展示了两个不同 $\beta$ 值下 $aV(an)$ 作为 $n$ 之函数的数值计算结果。实际的数值结果用符号表示，我们用虚线把它们连接起来以引导视线。由于数据来自数值模拟，它们带有（小的）统计误差，用水平误差棒表示。注意，两个图中坐标轴使用了不同的标度。垂直标度的选取使得两个图中静态势的形状相似。

这些图很好地展示了较大 $n$ 值处的线性上升行为，这是我们已从强耦合计算中推出的结论。对于短距离，数据点还显示出一种来自势的库仑部分的曲率。这在右侧的图（对应于较大的 $\beta$ 值）中更为明显。特别地，在小距离处，格点库仑势偏离连续形式。格点库仑势与格点自由玻色子传播子 [23] 有关，常被用来修正这种偏离 [24--26]。

\begin{figure}[htbp]
  \centering
  \includegraphics[width=0.9\textwidth]{fig35.png}
  \caption*{Fig. 3.5. \emph{用威尔逊规范作用量在两个不同耦合 $\beta$ 下计算的静态夸克势的数值数据。符号是我们为引导视线而连接起来的数值数据。虚线竖线画在与索末菲参数对应的距离处。数据取自 [21, 22]}.}
  \label{fig:35}
\end{figure}

\subsection{索末菲参数与格点间距}

我们已经预告，我们将引入索末菲参数作为确定格点间距 $a$ 的一种方式。在本书后面，我们将讨论确定 $a$ 的替代方法，例如利用强子的质量。尽管不同方法在给定的规范耦合下提供的 $a$ 值略有不同，但在连续极限下最终的物理结果应当一致。

索末菲参数 $r_0$ 是与静态势形状相关的一个特定距离。它的物理值是 $r_0 \approx 0.5$ fm。尽管我们只在下面讨论距离 $r_0$ 的精确定义，我们已经在图 3.5 的两个图中用虚线竖线标出了它的位置。因此，横轴上 $n = 0$ 与虚线之间的距离对应于 0.5 fm。注意，横轴上的变量 $n$ 是无量纲的——它不过是回路 $C$ 空间方向上格点间距的个数。因此，虚线竖线所标出的是格点单位下的索末菲标度，即 $r_0/a$。我们将在下面讨论如何从 $aV(an)$ 的数值数据中提取比值 $r_0/a$。

一旦算出索末菲标度 $r_0/a$，确定格点间距就很容易了。检视图 3.5 中的两个图时，我们发现原点（$n = 0$）与 $r_0/a$ 处虚线竖线之间的格点数对两个 $\beta$ 值是不同的。在图 3.5 的左侧图中，对应于 $\beta = 5.95$，我们发现虚线左侧约有 $n = 4.9$ 个格点。在 $\beta = 6.20$ 的图中，我们发现竖线左侧约有 $n = 7.4$ 个格点。由于两种情形中 $n = 0$ 与竖线之间的距离都对应于 0.5 fm，我们得到格点间距 $a$：
\begin{equation}
\begin{aligned}
  &a = 0.5\ \mathrm{fm}\,/\,4.9 \approx 0.102\ \mathrm{fm}
    && \text{for }\beta = 5.95 ,\\
  &a = 0.5\ \mathrm{fm}\,/\,7.4 \approx 0.067\ \mathrm{fm}
    && \text{for }\beta = 6.20 .
\end{aligned}
\end{equation}
我们最后需要讨论索末菲标度 $r_0$ 是如何定义的 [20]。索末菲标度并非直接基于势 $V(r)$，而是基于两个静态夸克之间的力 $F(r) = dV(r)/dr$。\footnote{$F$ 的这个定义与通常的定义相差一个负号。} 对于足够重的夸克，夸克-反夸克束缚态可以用有效的非相对论性薛定谔方程来描述，从而可以研究力 $F(r)$。通过与 $b\bar b$ 和 $c\bar c$ 谱的实验数据进行比较，人们发现
\begin{equation}
  F(r_0)\,r_0^2 = 1.65 \qquad \text{corresponds to} \quad r_0 \approx 0.5\ \mathrm{fm} .
\end{equation}
因此，我们需要从 $V(r)$ 的数值数据计算无量纲乘积 $F(r)\,r^2$，并确定这个乘积取值为 1.65 处的 $r = r_0$。为了便于理解这一确定过程，我们直接对 (3.62) 中参数化的势讨论索末菲参数。我们得到力
\begin{equation}
  F(r) = \frac{d}{dr}V(r)
  = \frac{d}{dr}\left(A + \frac{B}{r} + \sigma r\right)
  = -\frac{B}{r^2} + \sigma .
\end{equation}
于是，对于参数化的势，条件 (3.78) 为
\begin{equation}
  F(r_0)\,r_0^2 = -B + \sigma r_0^2 = 1.65 ,
\end{equation}
这意味着 $r_0 = \sqrt{(1.65 + B)/\sigma}$，或者用格点单位表示
\begin{equation}
  \frac{r_0}{a} = \sqrt{\frac{1.65 + B}{\sigma a^2}} .
\end{equation}
比值 $r_0/a$ 正是我们在图 3.5 中以及随后确定 $a$ 时所使用的量。数 $B$ 和 $\sigma a^2$ 可以通过把 $aV(an)$ 的数据拟合到（这是 $r = an$ 时的方程 (3.62)）\footnote{我们指出，这样的拟合已经假设势（至少在局部）具有 (3.62) 的形式。一种替代方法是直接在格点上离散化导数 $d/dr$ 来确定力 $F$（详见 [20]）。这种方法还有一个优点，即无需拟合 $\sigma a^2$ 和 $B$，而这样的拟合被大的 $n$ 值所主导，在大的 $n$ 处 $aV(an)$ 数据的信噪比变差。}
\begin{equation}
  aV(an) = Aa + \frac{B}{n} + \sigma a^2 n ,
\end{equation}
来确定，其中对不同的 $n$ 值拟合。于是，格点间距的确定可以总结如下：步骤 1：从 $aV(an)$ 的数值数据确定 $B$ 和 $\sigma a^2$。步骤 2：利用 (3.81) 计算无量纲数 $X = r_0/a$。步骤 3：格点间距由 $a = (0.5/X)\ \mathrm{fm}$ 给出。

在图 3.5 的讨论中，我们已经提到当增大 $\beta$ 时格点间距减小这一事实（另见 (3.77)）。在 [22] 中，格点间距 $a$ 是在几个 $\beta$ 值下为威尔逊作用量确定的，$a$ 对 $\beta$ 的依赖在 $5.7 \leq \beta \leq 6.92$ 范围内被参数化为
\begin{equation}
  a = r_0\,\exp\!\left[-1.6804 - 1.7331(\beta - 6) + 0.7849(\beta - 6)^2
    - 0.4428(\beta - 6)^3\right] .
\end{equation}
(3.83) 中参数化的形式受到重整化群的启发，我们现在简要讨论这一重要概念。

\subsection{重整化群与跑动耦合}

像规范耦合 $g$ 或夸克质量 $m$ 这样进入作用量泛函的耦合常数通常称为裸参数。这些并不是直接可观测的"物理"数值。只有通过计算诸如强子质量、弦张力或索末菲参数等可观测量，并把它们与实验值等同起来，人们才能以物理单位确定作用量裸参数的值。

格点作用量可以在各个方面有所不同。它们可以使用不同的导数离散化，或者格点网格（通常取为超立方）的结构可以变化。然而，当去除格点截断时，即在极限 $a \to 0$ 下，物理可观测量应当与实验值一致，并且变得与 $a$ 无关。一般来说，这将意味着裸参数对截断 $a$ 具有非平凡的依赖，即它们是函数 $g(a)$、$m(a)$ 等等。当我们令 $a \to 0$ 时，为了保持物理不变，裸参数的值必须改变。裸参数的这种跑动由所谓的重整化群来处理。为了简化对这一思想的讨论，我们考虑纯规范理论，其中只有一个裸耦合，即规范耦合 $g(a)$。

设 $P(g(a), a)$ 是一个在极限 $a \to 0$ 下取得其物理值 $P_0$ 的物理可观测量，
\begin{equation}
  \lim_{a\to 0} P(g(a), a) = P_0 .
\end{equation}
Callan 和 Symanzik 遵循 St\"uckelberg、Peterman、Gell-Mann 和 Low 早期对 QED 的建议，把物理不变的要求表述为一个微分方程：
\begin{equation}
  \frac{dP(g, a)}{d\ln a} = 0 \qquad \text{or, equivalently} \qquad
  \left(\frac{\partial}{\partial\ln a} + \frac{\partial g}{\partial\ln a}
    \frac{\partial}{\partial g}\right) P(g, a) = 0 .
\end{equation}
\begingroup\emergencystretch=4em 这个方程与改变标度的变换半群相联系，因此得名重整化群方程。（实际上，更精确地说，对于关联长度为 $\xi$ 的格点系统，这个方程的右端是\linebreak[0] $\order{(a/\xi)^2 \ln(a/\xi)}$。）第二项的系数函数称为 $\beta$ 函数（不要与我们的逆规范耦合混淆），\par\endgroup
\begin{equation}
  \beta(g) \equiv -\frac{\partial g}{\partial\ln a} ,
\end{equation}
它（除一个积分常数外）决定了耦合 $g$ 如何依赖于截断 $a$。

$\beta$ 函数可以在 $g = 0$ 附近展开为幂级数，其系数由微扰论确定。对于 SU($N$) 和 $n_f$ 个无质量夸克，结果为
\begin{equation}
\begin{aligned}
  \beta(g) &= -\beta_0\,g^3 - \beta_1\,g^5 + \order{g^7} ,\\
  \beta_0 &= \frac{1}{(4\pi)^2}\left(\frac{11}{3}N - \frac{2}{3}n_f\right) ,\\
  \beta_1 &= \frac{1}{(4\pi)^4}\left(\frac{34}{3}N^2 - \frac{10}{3}N\,n_f
    - \frac{N^2 - 1}{N}\,n_f\right) .
\end{aligned}
\end{equation}
对于纯规范 SU(3)，我们必须取 $N = 3$ 和 $n_f = 0$。展开的前两个系数是普适的，与正规化方案无关。然而，一般来说，$\beta$ 函数将取决于正规化的细节。

微分方程 (3.86) 连同 (3.87) 可以用分离变量法求解，得到
\begin{equation}
  a(g) = \frac{1}{\Lambda_L}\left(\beta_0 g^2\right)^{-\beta_1/(2\beta_0^2)}
    \exp\!\left(-\frac{1}{2\beta_0 g^2}\right)\left(1 + \order{g^2}\right) .
\end{equation}
积分常数 $\Lambda_L$ 用来通过在某个 $a$ 处固定 $g$ 的值来设定标度。反转关系 (3.88)，人们得到作为标度 $a$ 之函数的耦合 $g$，即所谓的跑动耦合，
\begin{equation}
  g(a)^{-2} = \beta_0 \ln\!\left(a^{-2}\Lambda_L^{-2}\right)
  + \frac{\beta_1}{\beta_0}\ln\!\left(\ln\!\left(a^{-2}\Lambda_L^{-2}\right)\right)
  + \order{1/\ln\!\left(a^2 \Lambda_L^2\right)} .
\end{equation}
因此，改变 $a$ 意味着 $g$ 相应改变，从而物理可观测量保持与设定标度的过程无关。$\Lambda_L$ 的值取决于正规化方案。不同的连续或格点作用量具有不同的 $\Lambda_L$ 值；然而，它们的比值可以由单圈微扰计算精确地联系起来。

事实上，微扰展开 (3.87) 的有效性必须加以检验。我们发现，当格点间距缩小时，跑动耦合也随之减小（对于 $n_f < 11N/2$）。格点间距趋于零对应于耦合 $g$ 趋于零，因为 $g = 0$ 确实是 $\beta$ 函数的一个零点。这一行为称为渐近自由。为诚实起见，我们必须指出，可能存在 $\order{\exp(-1/g^2)}$ 类型的非微扰贡献（对 $g > 0$ 非零，但不贡献于级数展开）。人们通常预期，这些贡献对足够小的 $g$ 不相关。此外，进行无穷体积极限与连续极限的顺序也是一个重要问题，参见 [27] 及其参考文献。

\subsection{真正的连续极限}

在第 1 章的结尾，我们已经简要提到如何实现格点场论的真正连续极限。然而，当时我们还不具备恰当地讨论这一极限的前提条件，并允诺稍后作详细讨论。对于纯规范理论的情形，我们现在可以提供这一讨论：由于我们分别在 (3.83) 和 (3.88) 中已经表明，格点间距 $a$ 随 $g$ 减小（$\beta$ 增大）而减小，我们得出结论，只需研究极限
\begin{equation}
  \beta \to \infty
\end{equation}
就能得到真正的连续极限 $a \to 0$。然而，在这一过程中有某些注意事项需要考虑。如果执行极限 (3.90)，那么我们研究 QCD 的盒子物理体积正比于 $a^4$，因而将收缩为零，除非我们也增加格点在空间方向（$N$ 个点）和时间方向（$N_T$ 个点）上的数目。在理想情况下，人们会首先执行所谓的热力学极限
\begin{equation}
  N \to \infty , \qquad N_T \to \infty ,
\end{equation}
并且只有在那一步之后才取连续极限 (3.90)。然而，由于在数值计算中这是不可行的，人们只能退而计算几个 $\beta$ 值下的物理可观测量，从而得到不同的 $a$ 值。格点数 $N$、$N_T$ 总是被选取为使得盒子的物理延伸
\begin{equation}
  L = a N , \qquad T = a N_T ,
\end{equation}
对不同 $a$ 值保持固定。在固定物理体积下研究结果的 $a$ 依赖性，使人们能够分析对标度 $a$ 的依赖，并把结果外推到 $a \to 0$。对 $a$ 依赖性的研究通常称为标度分析。然后可以对不同的物理尺寸 $L$、$T$ 重复向 $a = 0$ 的外推，这最终使人们能够把数据外推到无穷物理体积。

\section{其他规范群的格点规范理论}

尽管本书致力于格点 QCD，其规范群为 SU(3)，我们也简要讨论其他规范群的格点规范理论，特别是 SU($N$) 和 U(1)。SU(2) 之所以令人感兴趣，是因为它比 SU(3) 简单，而且是 SU(3) 的一个子群，因此被广泛研究。一些重要思想，例如拓扑荷和瞬子（即规范场经典运动方程的解），可以在 SU(2) 中表述。$N \geq 4$ 的群 SU($N$) 之所以令人感兴趣，是因为在连续理论中极限 $N \to \infty$ 是一个对解析研究很有吸引力的极限。因此，在格点上人们可以分析可观测量的 $N$ 依赖性。对于一般的 $N$，威尔逊作用量为
\begin{equation}
  S_G[U] = \frac{\beta}{N}\sum_{n\in\Lambda}\sum_{\mu<\nu}
    \Re\,\tr\!\left[1 - U_{\mu\nu}(n)\right] ,
\end{equation}
其中 $\beta = 2N/g^2$。基本方格 $U_{\mu\nu}(n)$ 是与 SU(3) 情形相同的有序乘积 (2.48)。我们迄今讨论的所有可观测量，特别是威尔逊圈和波利亚科夫圈，都可以不加改变地推广到一般的 $N$。

规范群 U(1) 之所以令人感兴趣，是因为它对应于 QED。作为一个阿贝尔群，它更容易实现，对这种情形的数值模拟是一个相当简单的练习，可以在 PC 上几小时内产生有趣的结果。与非阿贝尔群的主要区别在于，U(1) 的元素不是矩阵，而只是复相位，即我们可以把链接变量写成 $U_\mu(n) = \exp(\im A_\mu(n))$，其中 $A_\mu(n)$ 是实数。这意味着两个链接变量的对易子为零。因此，场强中的二次项为零，人们得到 $F_{\mu\nu}(x) = \partial_\mu A_\nu(x) - \partial_\nu A_\mu(x)$。不会出现自相互作用项。U(1) 的威尔逊作用量由下式给出
\begin{equation}
  S_G[U] = \beta\sum_{n\in\Lambda}\sum_{\mu<\nu}
    \Re\!\left(1 - U_{\mu\nu}(n)\right) .
\end{equation}
无需对色自由度取迹，基本方格只是复数的乘积，
$U_{\mu\nu}(n) = U_\mu(n)\,U_\nu(n+\hat\mu)\,U_\mu(n+\hat\nu)^*\,U_\nu(n)^*$
（星号表示复共轭）。耦合 $\beta$ 与 QED 的耦合 $e$ 通过 $\beta = 1/e^2$ 相联系。波利亚科夫圈和威尔逊圈可观测量可以再次使用，但无需取迹。

当把强耦合（小 $\beta$）展开应用于 U(1) 格点规范理论时，可以作出一个有趣的观察。人们仍然发现势中存在线性上升的项，这表明禁闭。另一方面，这个理论预期应当描述 QED，而 QED 并不显示出禁闭。这个谜团的解答是 U(1) 格点规范理论中的相变。在临界值 $\beta_{\mathrm{crit}} \approx 1.01$ 处，理论改变其行为，使得在 $\beta_{\mathrm{crit}}$ 以上，势的线性上升项消失，只剩下库仑势项。对于 SU($N$) 格点规范理论，人们普遍相信不会发生这样的相变，理论对所有 $\beta$ 都保持禁闭。然而，关于一种替代情形，参见 [28]。

\section*{参考文献}

\begin{enumerate}[label={[\arabic*]}]
  \item M.~Creutz: \emph{Quarks, Gluons and Lattices} (Cambridge University
        Press, Cambridge, New York 1983).
  \item H.~Georgi: \emph{Lie Algebras in Particle Physics}
        (Benjamin/Cummings, Reading, Massachusetts 1982).
  \item H.~F.~Jones: \emph{Groups, Representations and Physics} (Hilger,
        Bristol 1990).
  \item M.~Hamermesh: \emph{Group Theory and Its Application to Physical
        Problems} (Addison-Wesley, Reading, Massachusetts 1964).
  \item R.~Gilmore: \emph{Lie Groups, Lie Algebras, and Some of Their
        Applications}, Vol.~64 (John Wiley \& Sons, New York 1974).
  \item C.~Itzykson and J.-B.~Zuber: \emph{Quantum Field Theory} (McGraw-Hill,
        New York 1985).
  \item M.~E.~Peskin and D.~V.~Schroeder: \emph{An Introduction to Quantum
        Field Theory} (Addison-Wesley, Reading, Massachusetts 1995).
  \item S.~Weinberg: \emph{The Quantum Theory of Fields}, Vol.~1 and 2
        (Cambridge University Press, Cambridge, New York 1996).
  \item W.~H.~Press, S.~A.~Teukolsky, W.~T.~Vetterling, and B.~P.~Flannery:
        \emph{Numerical Recipes in C}, 2nd ed. (Cambridge University Press,
        Cambridge, New York 1999).
  \item L.~Giusti et al.: Int. J. Mod. Phys. A 16, 3487 (2001).
  \item S.~Elitzur: Phys. Rev. D 12, 3978 (1975).
  \item T.~DeGrand and C.~DeTar: \emph{Lattice Methods for Quantum
        Chromodynamics} (World Scientific, Singapore 2006).
  \item A.~M.~Polyakov: Phys. Lett. B 59, 82 (1975).
  \item C.~Borgs and E.~Seiler: Commun. Math. Phys. 91, 329 (1983).
  \item K.~G.~Wilson: Phys. Rev. D 10, 2445 (1974).
  \item K.~Osterwalder and E.~Seiler: Ann. Phys. 110, 440 (1978).
  \item J.~M.~Drouffe and J.~B.~Zuber: Phys. Rep. 102, 1 (1983).
  \item G.~Roepstorff: \emph{Path Integral Approach to Quantum Physics}
        (Springer, Berlin, Heidelberg, New York 1996).
  \item D.~H.~Perkins: \emph{Introduction to High Energy Physics}, 4th ed.
        (Cambridge University Press, Cambridge, New York 1999).
  \item R.~Sommer: Nucl. Phys. B 411, 839 (1994).
  \item M.~Guagnelli, R.~Sommer, and H.~Wittig: Nucl. Phys. B 535, 389 (1998).
  \item S.~Necco and R.~Sommer: Nucl. Phys. B 622, 328 (2002).
  \item C.~B.~Lang and C.~Rebbi: Phys. Lett. B 115, 137 (1982).
  \item C.~Michael: Nucl. Phys. B 259, 58 (1985).
  \item R.~G.~Edwards, U.~M.~Heller, and T.~R.~Klassen: Nucl. Phys. B 517, 377
        (1998).
  \item C.~R.~Allton et al.: Phys. Rev. D 65, 054502 (2002).
  \item E.~Seiler: in \emph{Applications of RG Methods in Mathematical
        Sciences} (RIMS, Kyoto University, Kyoto 2003).
  \item A.~Patrascioiu and E.~Seiler: Phys. Rev. Lett. 74, 1920 (1995).
\end{enumerate}
